\section{Introduction}

Income prediction is a practical classification problem with applications in labour economics, policy analysis, and benefit targeting. It is also a useful machine learning benchmark because errors do not all have the same implication: a model that predicts too many people as high-income behaves very differently from one that misses too many true high-income cases. That makes the task suitable for comparing not only overall accuracy, but also the type of mistakes made by different models.

In this project, we study the Adult income dataset from the UCI Machine Learning Repository~\cite{kohavi1996adult,uciadult}. Our objective is to determine whether a Random Forest classifier~\cite{breiman2001randomforest} or a Multi-Layer Perceptron (MLP)~\cite{rumelhart1986backprop,goodfellow2016deep} is better suited to predicting whether an individual's income exceeds \$50K. This is a meaningful comparison because the dataset is structured and mixed-type, which often suits tree ensembles, but it is still large enough to test whether a neural network can learn competitive decision boundaries after careful preprocessing.

Our study focuses on three practical questions. First, can the task be framed in a clean and reproducible pipeline from data preparation to final evaluation? Second, how differently do a tree ensemble and a neural network behave when trained on the same structured dataset? Third, do their computational costs tell a different story from their predictive scores? These questions shape the comparison carried out in the rest of the paper.

The rest of the paper is organised as follows. Section~\ref{sec:problem_setting} defines the task and dataset. Section~\ref{sec:methodology} presents the compared models. Section~\ref{sec:experimental_setup} describes the experimental protocol, and Section~\ref{sec:results} discusses the quantitative findings.
